% Pursuing Stacks (cotes 134-2 to 134-8): archive cote 134-2, site page 78 (PDF page 79).
% notes to ch. I. Draft compared: ps1.tex:31-78 ; ps1.tex:402-423.
% AI-transcribed by Claude Opus 5.5 (claude-opus-5-5, adaptive thinking, effort high), Message Batches API, batch ps1 (msgbatch_01Rv5hCaNTtUt4JHx1yqoHPL), 2026-11-12; not yet checked by eye.
\scanpage{134-2/78}
\mspage{62}

NOTES \hand{\uline{to Chapter I} \uline{and Appendix}.}

% NOTE: a large handwritten bracket in the left margin, with an arrow pointing right at its top, encloses the lines from "Chapter I" down to the note "These letters are reproduced ..."
\deleted{Chapter I}

(\over{\unsure{1}}{2}) \inserted[typed]{When making this} \xout{xxxmostxunexpected} suggestion about \inserted{there being} a "wind of disrepute for any foun%
\hand{-}dational matters whatever", I little suspected that the former friend \inserted{to whom} \deleted{whom}
I was \xout{\unsure{ask}} communicating my ponderings as they came, would take care of \xout{\illegible{1}}
providing a most unexpected confirmation. As a matter of fact, this let-%
ter never got an answer, nor was it even read! Upon my inquiry nearly one
year later, this colleague appeared \inserted[typed]{sincerely} surprised that \deleted{\inserted[typed]{\unsure{intensely}}} I could have expected
even for a minute that he might possibly read \deleted{this} \inserted{a} letter of mine on mathe-%
matical matters, well knowing the kind of ``general nonsense'' mathematics
\unsure{that} was to be expected from me ...

(\over{\unsure{2}}{1}) These letters are reproduced \inserted{as an ``appendix''} at the end of this chapter.

($^3$) For some particulars about a program of ``tame topology'', I refer to
``Esquisse d'un Programme'', sections 5 and 6, which is included in
Réflexions Mathématiques 1 .

($^4$) I have to apologize for this rash statement, as later correspondance
made me realize that ``Ronnie Brown and his friends'' do have stronger con-%
tact with ``geometry'' than I suspected, even though they are not too fami\hand{\unsure{-}}%
liar with algebraic geometry\hand{!}

($^5$) The ``Bangor group'' is made up by Ronnie Brown and Tim Porter as the
two fixed points, and a number of devoted research students. Moreover
Ronnie Brown is working in close contact with J.L. L\over{a}{o}day and J. Pradines\inserted{in France}\hand{.}

($^6$) Definitely only for \uline{strict} associativity.

($^7$) This idea is taken up again in section \over{\illegible{1}}{12}\hand{.} The statement made here
is a little rash, \deleted{\xout{andxthexsi} as well as} \inserted{as for} existence and uniqueness (in a
suitable sense) of this functor\hand{,} \deleted{\illegible{1}}. Compare note (1\over{\illegible{1}}{7}) below.

($^8$) Another important example is the structure of a \hand{``}torsor\hand{''} under a group
G (torsor = principal homogeneous space). When this group G is fixed,
the corresponding classifying topos $B_G$ is the natural purely algebraic sub%
\hand{-}stitute for the familiar \hand{``}classifying space\hand{''} for the discrete group G.

($^9$) Such a theory was developped in a seminar I gave at Buffalo in 1973.

($^{10}$) That \hand{$\frac{\unsure{i}}{2}$} is of order \over{\illegible{1}}{\unsure{1}} is heuristically clear, but will require a
proof none the less !

($^{11}$) (Added 23.2.83) I do'nt believe it now any more - and I do not \xout{rea}
really care - compare comments in section 11.

\margin{78}
